% AA-37 | Audyt I (2026-11-22); etykiety historyczne zachowane
\chapter{Algebraiczna niezależność i bazy transcendencji}
\label{aa:ch:transcendence}
\section{Motywacja: ile niezależnych parametrów ma rozszerzenie?}
Jeżeli $u,v$ są algebraicznie niezależnymi zmiennymi nad ciałem $K$,
to ciało $K(u,v)$ jest algebraiczne nad $K(u+v,uv)$:
oba elementy $u,v$ są pierwiastkami wielomianu
$T^2-(u+v)T+uv$. Ciała te nie muszą być równe.
Chcemy zatem liczyć nie liczbę \emph{wybranych} generatorów ciała,
lecz liczbę parametrów, których nie wiąże żadne niezerowe
równanie wielomianowe o współczynnikach z $K$.

\begin{definicja}\label{aa:def:alg-independent}
Niech $K\subseteq L$ będzie rozszerzeniem ciał, a $S\subseteq L$.
Zbiór $S$ nazywamy \emph{algebraicznie niezależnym nad $K$},
jeżeli dla każdego $m\ge1$, każdych parami różnych
$s_1,\ldots,s_m\in S$ i każdego niezerowego wielomianu
$P\in K[X_1,\ldots,X_m]$ zachodzi
$P(s_1,\ldots,s_m)\ne0$.
Zbiór pusty jest z definicji algebraicznie niezależny.
Zbiór $B\subseteq L$ jest \emph{bazą transcendencji rozszerzenia $L/K$},
jeżeli jest algebraicznie niezależny nad $K$, a każdy element $L$
jest algebraiczny nad ciałem $K(B)$ generowanym przez $K$ i $B$.
\end{definicja}

\begin{uwaga}
Niezależność zbioru jest własnością wszystkich jego
\emph{skończonych} podzbiorów. Z tego powodu definicja ma sens
także przy nieskończenie wielu parametrach; nie potrzebujemy
pierścienia wielomianów z nieskończonymi jednomianami.
Jeśli $B=\varnothing$ jest bazą transcendencji, oznacza to
po prostu, że $L/K$ jest algebraiczne.
\end{uwaga}

\begin{lemat}[Baza jako maksymalny zbiór niezależny]
\label{aa:lem:transcendence-maximal}
Niech $K\subseteq L$ będzie rozszerzeniem ciał i $B\subseteq L$.
Następujące warunki są równoważne:
\begin{enumerate}[label=(\roman*)]
 \item $B$ jest bazą transcendencji $L/K$;
 \item $B$ jest algebraicznie niezależny nad $K$ i nie istnieje
       algebraicznie niezależny $C$ taki, że $B\subsetneq C\subseteq L$.
\end{enumerate}
\end{lemat}
\begin{proof}
Załóżmy najpierw (i). Niech $x\in L\setminus B$.
Ponieważ $L/K(B)$ jest algebraiczne, istnieje moniczny
wielomian $m(T)=T^d+c_1T^{d-1}+\cdots+c_d\in K(B)[T]$
zerujący się w $x$. Każdy współczynnik $c_i$ jest ilorazem
dwóch wielomianów w skończenie wielu elementach $B$.
Po wybraniu wspólnego mianownika $q\ne0$ otrzymujemy
równanie $q x^d+q c_1x^{d-1}+\cdots+q c_d=0$,
którego współczynniki leżą w $K[b_1,\ldots,b_r]$
dla pewnych różnych $b_1,\ldots,b_r\in B$.
Wyrażenie $q$ jest niezerowe w $K(B)$, gdyż mianowniki
reprezentują niezerowe elementy ciała; niezależność $B$
gwarantuje ponadto, że jego zapis jako wielomian w $b_i$
nie jest zerowy. Otrzymujemy więc \emph{niezerowy} wielomian
w $r+1$ zmiennych zerujący się na $(b_1,\ldots,b_r,x)$.
To pokazuje zależność $B\cup\{x\}$. Dowolny właściwy
nadzbiór $B$ zawiera taki element $x$, a więc też jest zależny.

Załóżmy następnie (ii). Dla $x\in L\setminus B$ zbiór
$B\cup\{x\}$ jest zależny. Istnieją więc różne
$b_1,\ldots,b_r\in B$ i niezerowy
$P\in K[X_1,\ldots,X_r,T]$ z
$P(b_1,\ldots,b_r,x)=0$.
Wielomian $P$ musi mieć dodatni stopień względem $T$:
w przeciwnym razie naruszałby niezależność $B$.
Po podstawieniu $b_i$ dostajemy niezerowy wielomian
$P(b_1,\ldots,b_r,T)\in K(B)[T]$;
jego niezerowość wynika z faktu, że co najmniej jeden
niezerowy współczynnik przy potędze $T$ pozostaje niezerowy
po podstawieniu niezależnych $b_i$.
W konsekwencji $x$ jest algebraiczny nad $K(B)$.
Elementy $x\in B$ oczywiście należą do $K(B)$.
Wobec tego każde $x\in L$ jest algebraiczne nad $K(B)$,
co daje (i).
\end{proof}

\begin{twierdzenie}[Istnienie bazy transcendencji]
\label{aa:tw:transcendence-existence}
Każde rozszerzenie ciał $K\subseteq L$ posiada bazę transcendencji.
\end{twierdzenie}
\begin{proof}
Niech $\mathcal I$ będzie rodziną algebraicznie niezależnych
podzbiorów $L$, uporządkowaną przez inkluzję.
Rodzina nie jest pusta, gdyż $\varnothing\in\mathcal I$.
Jeżeli $\mathcal C\subseteq\mathcal I$ jest łańcuchem,
to $U=\bigcup_{C\in\mathcal C}C$ jest jego ograniczeniem górnym.
Aby sprawdzić niezależność $U$, weźmy dowolne
$s_1,\ldots,s_m\in U$. Dla każdego $s_i$ istnieje
$C_i\in\mathcal C$ zawierający $s_i$. Ponieważ $\mathcal C$
jest łańcuchem i lista $C_1,\ldots,C_m$ jest skończona,
jedno z $C_i$ zawiera wszystkie pozostałe. Zatem
$s_1,\ldots,s_m$ leżą razem w pewnym zbiorze niezależnym,
więc nie spełniają niezerowej relacji wielomianowej.
Lemat Zorna, wymagający w ogólności aksjomatu wyboru,
daje maksymalny $B\in\mathcal I$.
Na mocy Lematu~\ref{aa:lem:transcendence-maximal}
jest to baza transcendencji.
\end{proof}

\section{Wymiana elementów}
\begin{lemat}[Wymiana algebraiczna]
\label{aa:lem:transcendence-exchange}
Niech $K\subseteq L$ będzie rozszerzeniem ciał,
$B\subseteq L$ zbiorem algebraicznie niezależnym nad $K$,
a $b\in B$. Jeśli $x\in L$ jest algebraiczny nad $K(B)$,
lecz transcendentalny nad $F=K(B\setminus\{b\})$,
to $b$ jest algebraiczny nad $F(x)$.
\end{lemat}
\begin{proof}
Element $b$ jest transcendentalny nad $F$, ponieważ $B$ jest
niezależny. Z algebraiczności $x$ nad $F(b)$ wynika istnienie
niezerowego wielomianu $P(T)\in F(b)[T]$ zerującego się w $x$.
Usuwając mianowniki współczynników, otrzymujemy
niezerowy wielomian $Q(U,V)\in F[U,V]$
taki, że $Q(b,x)=0$; wielomian pozostaje niezerowy,
ponieważ podstawienie formalnego $U=b$ jest iniektywne
na $F[U]$.
Zapiszmy $Q(U,V)=\sum_{j=0}^r q_j(V)U^j$,
gdzie $q_j\in F[V]$ i co najmniej jeden $q_j$ jest niezerowy.
Ponieważ $x$ jest transcendentalny nad $F$,
niezerowy $q_j$ nie może spełniać $q_j(x)=0$.
Zatem $Q(U,x)\in F(x)[U]$ jest niezerowy i zeruje się
dla $U=b$. To jest równanie algebraiczne elementu $b$
nad $F(x)$.
\end{proof}

\begin{lemat}[Nierozkładalność po dołączeniu parametru]
\label{aa:lem:transcendence-irreducible-pure}
Niech $F$ będzie ciałem, $z$ elementem transcendentalnym nad $F$,
a $h(T)\in F[T]$ wielomianem nierozkładalnym dodatniego stopnia.
Wówczas $h$ pozostaje nierozkładalny w $F(z)[T]$.
\end{lemat}
\begin{proof}
Przypuśćmy, że $h=uv$ w $F(z)[T]$,
gdzie $u,v$ mają dodatnie stopnie względem $T$.
Pierścień $F[z]$ jest pierścieniem wielomianów nad ciałem,
a więc dziedziną z jednoznacznością rozkładu.
Wielomian $h$ jest prymitywny jako element $F[z][T]$:
jego niezerowe współczynniki należą do $F$ i ich
największy wspólny dzielnik jest jednostką.
Z lematu Gaussa wynika istnienie rozkładu
$h=\widetilde u\widetilde v$ w $F[z][T]$,
gdzie oba czynniki mają dodatni stopień względem $T$.
Traktując tę równość jako tożsamość w $F[T][z]$, mamy
$\deg_z h=0$. Ponieważ $F[T]$ jest dziedziną,
$\deg_z(\widetilde u\widetilde v)=
\deg_z\widetilde u+\deg_z\widetilde v$.
Oba stopnie muszą zatem wynosić $0$, czyli
$\widetilde u,\widetilde v\in F[T]$.
Otrzymujemy rozkład $h$ nad $F$, wbrew założeniu.
\end{proof}

\begin{lemat}[Zmienna istotna w funkcji wymiernej]
\label{aa:lem:transcendence-essential-variable}
Niech $D$ będzie zbiorem algebraicznie niezależnym nad ciałem $K$,
$A\subseteq D$, a $r\in K(D)\setminus K(A)$.
Istnieje $d\in D\setminus A$ takie, że
$r\notin K(D\setminus\{d\})$.
\end{lemat}
\begin{proof}
Z definicji ciała generowanego przez $D$ element $r$ jest
ilorazem dwóch wielomianów zależnych od skończenie wielu
elementów $D$. Niech $E\subseteq D\setminus A$ będzie
skończonym zbiorem obejmującym wszystkie elementy spoza $A$
użyte w takim ilorazie. Wtedy $r\in K(A)(E)$.
Wykażemy najpierw własność przecięcia: dla każdego $d\in E$
\[
 K(A)(E)\cap K(D\setminus\{d\})
   =K(A)(E\setminus\{d\}).
\]
Inkluzja od prawej do lewej wynika z definicji.
Dla inkluzji odwrotnej weźmy $u$ należące do przecięcia.
Jako funkcję wymierną zmiennej $d$ możemy je zapisać
$u=P(d)/Q(d)$, gdzie
$P,Q\in K(A)(E\setminus\{d\})[T]$, $Q\ne0$.
Jednocześnie $u\in K(D\setminus\{d\})$.
Ponieważ $d$ jest transcendentalny nad tym ostatnim ciałem,
z równości $P(d)=uQ(d)$ wynika równość każdego
współczynnika tych dwóch wielomianów.
Wybierzmy niezerowy współczynnik $q_j$ wielomianu $Q$.
Wówczas odpowiadający mu współczynnik $p_j$ wielomianu $P$
spełnia $p_j=uq_j$, skąd
$u=p_j/q_j\in K(A)(E\setminus\{d\})$.

Gdyby dla każdego $d\in E$ zachodziło
$r\in K(D\setminus\{d\})$, to powyższa własność pozwalałaby
usuwać kolejno wszystkie elementy skończonego zbioru $E$
z przedstawienia $r$.
Ostatecznie otrzymalibyśmy $r\in K(A)$,
sprzecznie z założeniem. Istnieje zatem wymagane $d$.
\end{proof}

\begin{twierdzenie}[Jednoznaczność liczby elementów bazy]
\label{aa:tw:transcendence-degree}
Niech $K\subseteq L$ będzie dowolnym rozszerzeniem ciał.
Jeżeli $B,C\subseteq L$ są bazami transcendencji $L/K$,
to mają tę samą moc. Moc tę nazywamy
\emph{stopniem transcendencji} i oznaczamy $\operatorname{trdeg}_K L$.
\end{twierdzenie}
\begin{proof}
\emph{Krok 1: zamiana jednego elementu w bazie.}
Niech $D$ będzie bazą transcendencji $L/K$, a $A\subseteq D$
podrodziną, której elementy chcemy zachować. Weźmy
$x\in L$ transcendentalny nad $K(A)$.
Ponieważ $L/K(D)$ jest algebraiczne, element $x$ jest
algebraiczny nad $K(D)$. Jego moniczny wielomian minimalny
$m_x(T)\in K(D)[T]$ nie może mieć wszystkich współczynników
w $K(A)$, gdyż wtedy $x$ byłby algebraiczny nad $K(A)$.
Ponieważ co najmniej jeden współczynnik $m_x$ nie należy
do $K(A)$, Lemat~\ref{aa:lem:transcendence-essential-variable}
pozwala wybrać $d\in D\setminus A$ tak, aby ten
współczynnik nie należał do $F=K(D\setminus\{d\})$.

Gdyby $x$ był algebraiczny nad $F$, jego moniczny
wielomian minimalny $h(T)\in F[T]$ pozostałby nierozkładalny
po dołączeniu transcendentalnego nad $F$ elementu $d$,
na mocy Lematu~\ref{aa:lem:transcendence-irreducible-pure}.
Ponieważ $F(d)=K(D)$, mielibyśmy $m_x=h$, co przeczyłoby
wyborowi $d$. Stąd $x$ jest transcendentalny nad $F$.
Z Lematu~\ref{aa:lem:transcendence-exchange} wynika,
że $d$ jest algebraiczny nad $F(x)$.
Połóżmy $D'=(D\setminus\{d\})\cup\{x\}$.
Niezależność $D'$ wynika z transcendentalności $x$ nad $F$,
a fakt, że $L/K(D')$ jest algebraiczne, z dwóch zdań:
$L$ jest algebraiczne nad $K(D)=F(d)$ oraz $d$ jest algebraiczny
nad $F(x)=K(D')$. Przechodniość algebraiczności kończy dowód,
że $D'$ jest bazą zawierającą $A\cup\{x\}$.

\emph{Krok 2: bazy skończone.}
Załóżmy, że $C$ ma $m<\infty$ elementów.
Dla dowolnych $m+1$ różnych elementów $b_1,\ldots,b_{m+1}\in B$
próbujemy kolejno wprowadzić je do bazy $C$.
Po wprowadzeniu $b_1,\ldots,b_j$ aktualna baza zawiera
$A_j=\{b_1,\ldots,b_j\}$; kolejny element $b_{j+1}$ jest
transcendentalny nad $K(A_j)$, ponieważ jest częścią
algebraicznie niezależnego zbioru $B$.
Z Kroku 1 można więc wymienić element spoza $A_j$ na
$b_{j+1}$, nie zmieniając liczby elementów bazy.
Po $m$ krokach baza ma nadal $m$ elementów, wszystkie należące
do $B$. Element $b_{m+1}$ jest nad nimi transcendentalny,
a więc nie może należeć do algebraicznego rozszerzenia
$K(C')\subseteq L$ wyznaczonego przez tę bazę $C'$.
Sprzeczność dowodzi $|B|\le m$.
W szczególności $B$ jest skończona; zamieniając role $B,C$,
uzyskujemy $|C|\le|B|$, a więc równość.

\emph{Krok 3: bazy nieskończone.}
Jeżeli żadna baza nie jest skończona, to dla każdego $c\in C$
istnieje skończony zbiór $B_c\subseteq B$ taki, że
$c$ jest algebraiczny nad $K(B_c)$: algebraiczność $c$
nad $K(B)$ jest wyrażona równaniem o skończonej liczbie
współczynników, z których każdy zawiera tylko skończenie
wiele elementów $B$.
Niech $B_*=\bigcup_{c\in C}B_c$.
Wtedy każdy $c\in C$ jest algebraiczny nad $K(B_*)$.
Każdy element $L$ jest z kolei algebraiczny nad $K(C)$,
więc, na mocy przechodniości algebraiczności,
$L$ jest algebraiczne nad $K(B_*)$.
Zbiór $B_*$ jest niezależny, bo $B_*\subseteq B$,
a zatem jest bazą transcendencji.
Jeśli $B_*\subsetneq B$, wybór $b\in B\setminus B_*$ dawałby
element transcendentalny nad $K(B_*)$, sprzecznie
z algebraicznością $L/K(B_*)$. Stąd $B=B_*$ i
\[
 |B|\le\sum_{c\in C}|B_c|\le |C|\cdot\aleph_0=|C|,
\]
gdzie ostatnia równość wynika z arytmetyki nieskończonych
liczb kardynalnych przy aksjomacie wyboru.
Ściślej, wybieramy dla każdego $c\in C$ skończony zbiór
$B_c$; ponieważ $C$ jest nieskończony, zbiór par
$(c,b)$ z $c\in C$ i $b\in B_c$ wstrzykuje się
w $C\times\mathbb N$ po wybraniu numeracji każdego
$B_c$. Stąd $|B|\le |C\times\mathbb N|=|C|$.
Po zamianie ról $B,C$ otrzymujemy odwrotną nierówność,
więc $|B|=|C|$.
\end{proof}

\begin{twierdzenie}[Dodawanie stopni transcendencji w wieży ciał]
\label{aa:tw:transcendence-tower}
Niech $K\subseteq E\subseteq L$ będzie wieżą rozszerzeń ciał.
Jeżeli $U\subseteq E$ jest bazą transcendencji $E/K$,
a $V\subseteq L$ jest bazą transcendencji $L/E$,
to $U\cup V$ jest bazą transcendencji $L/K$. W szczególności
\[
 \operatorname{trdeg}_K L
 =\operatorname{trdeg}_K E+\operatorname{trdeg}_E L,
\]
gdzie dodawanie mocy zbiorów rozumiemy także dla baz
nieskończonych.
\end{twierdzenie}
\begin{proof}
\emph{Krok 1: rozłączność i niezależność.}
Element $v\in V$ nie może należeć do $E$: wielomian
$T-v\in E[T]$ przeczyłby niezależności $V$ nad $E$.
Zatem $U\cap V=\varnothing$.
Przypuśćmy, że różne $u_1,\ldots,u_r\in U$ i
$v_1,\ldots,v_s\in V$ spełniają niezerową relację
$P(u_1,\ldots,u_r,v_1,\ldots,v_s)=0$, gdzie
$P\in K[X_1,\ldots,X_r,Y_1,\ldots,Y_s]$.
Zapiszmy $P=\sum_\alpha P_\alpha(X_1,\ldots,X_r)Y^\alpha$.
Ponieważ $P\ne0$, co najmniej jeden $P_\alpha$ jest
niezerowy. Niezależność $U$ nad $K$ gwarantuje
$P_\alpha(u_1,\ldots,u_r)\ne0$.
Po podstawieniu $u_i$ otrzymujemy więc niezerowy
wielomian o współczynnikach w $K(U)\subseteq E$,
który zeruje się na $(v_1,\ldots,v_s)$; sprzeczność
z niezależnością $V$ nad $E$.
Jeżeli $r=0$ lub $s=0$, argument sprowadza się
bezpośrednio do odpowiedniej niezależności $U$ lub $V$.

\emph{Krok 2: algebraiczność nad ciałem parametrów.}
Z definicji bazy transcendencji każdy $e\in E$ jest
algebraiczny nad $K(U)$, a więc również nad większym
ciałem $K(U\cup V)$. Oznacza to, że rozszerzenie
$E(V)/K(U\cup V)$ jest algebraiczne: jego elementy
powstają z elementów $E$ i $V$ przez skończenie wiele
działań ciała, zaś elementy algebraiczne nad danym
ciałem tworzą podciało.
Ponieważ $L/E(V)$ jest algebraiczne, z przechodniości
algebraiczności wynika, że $L/K(U\cup V)$ również
jest algebraiczne. Krok 1 i definicja bazy kończą dowód.
Równość stopni otrzymujemy z jednoznaczności mocy baz
(Twierdzenie~\ref{aa:tw:transcendence-degree}) i rozłączności
$U,V$.
\end{proof}

\begin{przyklad}[Wieża ciał i rozszerzenie algebraiczne]
\label{aa:prz:transcendence-tower-rational}
Niech $u,v$ będą algebraicznie niezależne nad $K$,
a $L=K(u,v)$ i $E=K(u+v)$. Podstawienie
$S\mapsto u+v$ jest iniekcją $K[S]\to K[u,v]$,
gdyż niezerowy wielomian w $u+v$ ma niezerowy
współczynnik przy najwyższej potędze $u$.
Stąd $\operatorname{trdeg}_K E=1$.
Ponieważ $u=(u+v)-v$, otrzymujemy $L=E(v)$.
Element $v$ jest transcendentalny nad $E$. Istotnie,
po oczyszczeniu mianowników ewentualnej relacji
nad $E=K(u+v)$ dostalibyśmy niezerowy wielomian
$P(S,T)\in K[S,T]$ z $P(u+v,v)=0$.
Jednak podstawienie $u\mapsto u+v$, $v\mapsto v$
jest automorfizmem $K[u,v]$ z odwrotnością
$u\mapsto u-v$, $v\mapsto v$; niezerowy $P$
nie może zatem zerować się na $(u+v,v)$.
Wobec tego $\operatorname{trdeg}_E L=1$, co daje $2=1+1$.

Dla porównania rozszerzenie $K(u,v)/K(u+v,uv)$
jest algebraiczne: obie zmienne spełniają równanie
$T^2-(u+v)T+uv=0$. Dodanie rozszerzenia algebraicznego
nie zmienia stopnia transcendencji, co jest przypadkiem
twierdzenia dla $\operatorname{trdeg}_E L=0$.
\end{przyklad}

\begin{przyklad}[Nieskończona baza transcendencji]
\label{aa:prz:transcendence-infinite}
Niech $(t_j)_{j\ge1}$ będzie ciągiem elementów algebraicznie
niezależnych nad ciałem $K$ i niech
$L=K(t_1,t_2,\ldots)$. Zbiór $B=\{t_j:j\ge1\}$
jest bazą transcendencji $L/K$: jest niezależny z założenia,
a każdy element $L$ należy już do $K(B)$.
Ponadto baza ta jest przeliczalnie nieskończona,
więc $\operatorname{trdeg}_K L=\aleph_0$.

Rozważmy również
$C=\{t_1+t_2,t_2,t_3,t_4,\ldots\}$.
Zachodzi $K(C)=L$, bo $t_1=(t_1+t_2)-t_2$.
Aby wykazać niezależność $C$, załóżmy, że pewien niezerowy
wielomian $P\in K[X_1,\ldots,X_m]$ zeruje się po
podstawieniu skończenie wielu różnych elementów $C$.
Wówczas otrzymujemy niezerową relację wielomianową między
skończenie wieloma $t_j$, gdyż podstawienie
$t_1\mapsto t_1+t_2$ i $t_j\mapsto t_j$ dla $j\ge2$
jest automorfizmem pierścienia
$K[t_1,t_2,\ldots]$, o odwrotności
$t_1\mapsto t_1-t_2$ oraz $t_j\mapsto t_j$ dla $j\ge2$.
To przeczy niezależności $B$. Zatem $C$ również jest
bazą transcendencji $L/K$.
Przykład pokazuje też, że nawet dla bazy nieskończonej
sprawdzanie niezależności sprowadza się do relacji
obejmujących tylko skończenie wiele elementów.
\end{przyklad}

\section{Przykłady i kontrprzykłady}
\begin{przyklad}[Dwie różne bazy w tym samym ciele]
\label{aa:prz:transcendence-rational}
Niech $u,v$ będą niezależnymi zmiennymi nad ciałem $K$.
Z definicji $\{u,v\}$ jest bazą transcendencji $K(u,v)/K$.
Połóżmy $s=u+v$ i $p=uv$.
Na mocy jednoznaczności zapisu wielomianu symetrycznego
(Twierdzenie~\ref{aa:tw:fundamental-symmetric}) $s,p$
są algebraicznie niezależne nad $K$.
Ponadto zarówno $u$, jak i $v$ spełniają
$T^2-sT+p=0$ nad $K(s,p)$.
Ciało $K(u,v)$ jest więc algebraiczne nad $K(s,p)$;
zatem $\{s,p\}$ jest drugą bazą transcendencji i
$\operatorname{trdeg}_K K(u,v)=2$.
W charakterystyce $2$ rachunek działa tak samo: współczynnik
$-s$ jest równy $s$, lecz tożsamość iloczynu
$(T-u)(T-v)=T^2-sT+p$ pozostaje prawdziwa.
\end{przyklad}

\begin{przyklad}[Dwie bazy jednoelementowe]
Niech $t$ będzie transcendentalny nad $K$ i $m\ge2$.
Zbiór $\{t^m\}$ jest niezależny, ponieważ
$P(t^m)=0$ dla niezerowego $P\in K[T]$ dawałoby
niezerową relację wielomianową dla $t$.
Element $t$ jest algebraiczny nad $K(t^m)$:
spełnia moniczne równanie $T^m-t^m=0$.
Dlatego $\{t\}$ i $\{t^m\}$ są różnymi bazami
transcendencji tego samego rozszerzenia $K(t)/K$.
\end{przyklad}


\begin{przyklad}[Baza transcendencji a nierozdzielczość]
\label{aa:prz:transcendence-inseparable}
Niech $K$ będzie ciałem charakterystyki $p>0$,
a $t$ elementem transcendentalnym nad $K$.
Ciało $K(t)$ jest algebraiczne nad $K(t^p)$, gdyż
$t$ spełnia $T^p-t^p=0$. Zatem zbiory $\{t\}$
i $\{t^p\}$ są bazami transcendencji $K(t)/K$.
Ponadto $[K(t):K(t^p)]=p$. Istotnie, gdyby
$1,t,\ldots,t^{p-1}$ były liniowo zależne nad $K(t^p)$,
to po oczyszczeniu mianowników mielibyśmy
$\sum_{i=0}^{p-1}A_i(t^p)t^i=0$ z pewnymi
$A_i\in K[T]$, nie wszystkimi zerowymi.
W rozwinięciu tej sumy potęgi $t$ pochodzące
od różnych indeksów $i$ mają różne reszty modulo $p$;
nie mogą się więc skracać. Otrzymujemy sprzeczność
z transcendentalnością $t$. Natomiast równanie
$T^p-t^p=0$ dowodzi nierówności odwrotnej dla stopnia.
Jest to rozszerzenie czysto nierozdzielcze: wielomian
minimalny $T^p-t^p$ ma pochodną formalną równą zeru.
Zależność algebraiczna potrzebna
w definicji bazy nie wymaga rozdzielczości;
w szczególności twierdzenie o jednoznaczności stopnia
transcendencji działa również w dodatniej charakterystyce.
\end{przyklad}

\begin{przyklad}[Kontrprzykład: dwa generatory nie muszą być niezależne]
\label{aa:prz:transcendence-dependent}
W ciele $K(t)$ zbiór $\{t,t^2\}$ generuje to samo
rozszerzenie co $\{t\}$, ale nie jest niezależny:
niezerowy wielomian $X_2-X_1^2\in K[X_1,X_2]$
zeruje się na parze $(t,t^2)$.
Liczba generatorów ciała nie musi być zatem równa
stopniowi transcendencji.
\end{przyklad}


\begin{przyklad}[Jawna wymiana parametru w bazie]
\label{aa:prz:transcendence-exchange-explicit}
Niech $u,v$ będą algebraicznie niezależne nad ciałem $K$,
$L=K(u,v)$ oraz $x=u+v^2$. Zbiory $\{u,v\}$ i
$\{x,v\}$ są bazami transcendencji $L/K$.
Rzeczywiście, dla niezerowego $P(S,T)\in K[S,T]$
założenie $P(x,v)=0$ oznaczałoby
$P(u+v^2,v)=0$. Podstawienie $S\mapsto U+V^2$,
$T\mapsto V$ jest automorfizmem $K[S,T]$ z odwrotnością
$U\mapsto S-T^2$, $V\mapsto T$; dlatego
$P(U+V^2,V)$ pozostaje niezerowym wielomianem.
Algebraiczna niezależność $u,v$ wyklucza jego zerowanie.
Ponadto $u=x-v^2\in K(x,v)$, więc $K(x,v)=L$.
To pokazuje bezpośrednio, że wymiana elementu bazy
może zachować nawet równość generowanych ciał,
a nie tylko algebraiczność ich rozszerzenia.
\end{przyklad}
